\documentclass[10pt,a4paper]{amsart}
\usepackage[margin=19mm]{geometry}
\usepackage{amsmath,amssymb,mathtools}
\usepackage[scr=rsfs,cal=euler]{mathalfa}
\usepackage{xfrac}
\usepackage[unicode,hidelinks,bookmarks=false]{hyperref}
\newcommand{\ie}{i.e.\ }
\newcommand{\cf}{cf.\ }
\newcommand{\resp}{respectively\ }
\newcommand{\Subsection}{\subsection}
\providecommand{\nobreakdash}{\nobreak\hbox{-}}
\setlength{\emergencystretch}{2em}
\multlinegap0pt
\usepackage{amssymb}
\usepackage[shortlabels]{enumitem}
\newtheorem{theorem}{Theorem}
\title{On the boundary behavior of analytic functions in the unit disc}
\author{Spyros Pasias}
\date{Source: arXiv:2305.10870v4 (CC BY 4.0)}
\begin{document}
\maketitle

\noindent\emph{Source and licence.} On the boundary behavior of analytic functions in the unit disc. Version arXiv:2305.10870v4 (CC BY 4.0), distributed by arXiv under the Creative Commons Attribution 4.0 International licence, http://creativecommons.org/licenses/by/4.0/, as recorded in the arXiv licence metadata for this exact version and shown on the abstract page https://arxiv.org/abs/2305.10870v4. Full licence notice: \texttt{licenses/arxiv-2305.10870-CC-BY-4.0.txt}.\par\smallskip

\noindent\textit{Editorial scope.} Counted unit: the article's theorem theorem (source0.tex lines 109--114). The article's complete original proof of it is delivered with it (source0.tex lines 116--124, one \texttt{proof} environment of 9 lines). No mathematics is written, repaired or re-derived: the statement and the proof are the article's own lines, in the article's own notation, and no retained line is substituted or re-flowed.  No further source-local context is needed: every cross-reference of every retained line resolves inside the two delivered spans.  Nothing was omitted from any retained span, so the package carries no omission record.  No retained line contains a citation, so the wrapper carries no bibliography and no bibliography entry.  No retained line contains a figure environment, an image-inclusion command or drawing code, so no image is delivered and the package declares no asset dependency.  Editorial formatting only, in addition: an AMS \texttt{amsart} preamble that restates the article's own declarations and macros used by the retained lines, namely the environments \texttt{theorem} on separate counters, together with the standard packages those lines need.  The retained environments print in this document as Theorem 1 in order of appearance, whereas the article prints the same items with the numbering of its own full document.  The complete native source of the article as distributed in the arXiv e-print for this exact version is bundled unchanged as \texttt{sources/141693/source0.tex}.
\begin{theorem} Let $E_1\subset E_2$ be subsets of the unit circle $T$ and suppose that the measure of $E_2$ is $2\pi$. The following conditions are necessary and sufficient for the existence of an $f\in H^\infty$ such that the radial limits of $f$ fail to exist precisely at $E_1$ and it has vanishing unrestricted limits precisely at $T\setminus E_2$:
\begin{enumerate}
\item $E_1$ is of type $G_{\delta\sigma}$ and of measure zero.
\item $E_2$ is of type $F_\sigma$.
\end{enumerate}
\end{theorem}

\begin{proof}
\textbf{(Necessity)} It is well known that the exceptional set of points in $T$ such that $f$ fails to have radial
 limits is of type $G_{\delta\sigma}$ and by {Fatou's Theorem} is of measure zero. 
 Also if $f$ has unrestricted limits on $T\setminus E_2$ then $E_2=\bigcup_{n=1}^{\infty} F_n$ where $F_n$ are the closed sets consisting 
 of all points in $E_2$ any neighborhood of which contains points $t_1,t_2\in D$ such that $|f(t_1)-f(t_2)|>\frac{1}{n}$. Since clearly each $F_n$ is closed it follows that $E_2$ is of type $F_{\sigma}$ and the proof of necessity is complete.\\
\textbf{(Sufficiency)} Since $E_1$ is of type $G_{\delta\sigma}$ and of measure zero by {Kolesnikov's Theorem} there exists a bounded analytic function $f_1$ that fails to have radial limits exactly on $E_1$. By adding a suitable constant we may assume that $f_1$ does not have vanishing radial limits anywhere on $T$.\\
 Now since $E_2$ is $F_\sigma$ it follows $T\setminus E_2$ is of type $G_\delta$ and therefore by {Theorem 4} there exists a non zero bounded analytic function $f_2$ such that $f_2=0$ exactly on $T\setminus E_2$ and such that the radial limits of $f_2$ exist at all points of $T$.
The function $F\equiv f_1f_2$ satisfies the conditions of the theorem hence the proof is complete. 
\end{proof}

\end{document}
